% r03_r1_setting: the cumulant bound for the step Gaussian, part I: the setting, the lattice facts, the constants, and the kernel and decoupling lemmas. Input from sec/r03_fill_K5a.tex before Theorem rot:thm:R1.

\subsubsection*{The cumulant bound for the step Gaussian}

The fluctuation integral of every block-spin step of both flows has the form $\int\chi_b\,e^{V}d\mu_P$: $\mu_P$ is a centred Gaussian whose precision is coercive and has exponentially decaying off-diagonal entries, $V$ is a real polynomial with exponentially decaying coefficients, and $\chi_b$ is a product over sites of indicators of a fixed convex body scaled by the threshold $b$. \Cref{rot:thm:R1} below is the order-$\bar n$ cumulant expansion of such an integral, with an error proportional to the volume of the support of $V$ and, in the lower bound, a normalisation term proportional to $|\Lambda|$. The printed cluster-expansion lemma of Benfatto and coauthors, which King cites for this step, is stated for the nearest-neighbour Markov field with a cutoff growing with the distance to the support of the interaction; its proof uses the Markov property and the nearest-neighbour formulas for conditional Gaussians in four places. The proof below keeps its architecture (cells with corridors, four shifted pavements, the cutoff moved from the corridors into the cells) and replaces those four uses by the decoupling and conditioning lemmas of this part, which hold for any coercive precision with exponentially decaying entries.

\paragraph{The setting.} Let $\mathbb{L}$ be one of: the lattice $\Z^3$; a cubic torus $(\Z/N\Z)^3$; a skew torus $\Z^3/mG_\pm$, $m\in\N$, where $G_+\defeq\langle(40,-30,0),(30,40,0),(0,0,50)\rangle$ and $G_-$ is its mirror image under $x_2\mapsto-x_2$. These are the frame period lattices of the skew tori of \Cref{rot:prop:SK} in the units of a step lattice: $\Gamma'=R^{\mathrm T}(50MJ\,\Z^3)=MJ\,G_+$ for $R\DJ$ in $Y$'s frame and $RP_J=MJ\,G_-$ for $R\DJ$ in $X$'s frame, so $m=MJ$ divided by the spacing of the step, an integer because every block lattice contains $P_J$ (\Cref{rot:lem:adm}(iii)). Distances are quotient Euclidean distances, $|x-y|\defeq\min_g|\tilde x-\tilde y-g|$ over the period lattice (a metric), and $\dist(x,S)\defeq\min_{y\in S}|x-y|$. $\Lambda\subset\mathbb{L}$ is finite, $z=(z_x)_{x\in\Lambda}$ with $z_x\in\R^n$, and a complex Hermitian operator on $\C^\Lambda$ is read in its real form (blocks $\bigl(\begin{smallmatrix}\Re&-\Im\\\Im&\Re\end{smallmatrix}\bigr)$), which preserves the spectrum and the block operator norms. The hypotheses are:
\begin{enumerate}[label=\textup{(R\arabic*)},leftmargin=*]
\item $\mu_P\defeq N(0,P^{-1})$ with $P$ real symmetric on $(\R^n)^\Lambda$ and $P\ge\gamma_0>0$ (no upper bound is needed);
\item $\|P(x,y)\|\le c_Pe^{-\delta_P|x-y|}$ for $x\ne y$ (block operator norms); the diagonal blocks are arbitrary;
\item $V=\sum_{Y\subset J}V_Y$ with $J\subset\Lambda$, where $V_Y$ is a real polynomial in $(z_y)_{y\in Y}$ of degree at most $D$, without constant term, each of whose monomials involves every site of $Y$ (its \emph{support}); $\ell(Y)$ is the length of a shortest connected graph on a vertex set containing $Y$ (edges between points of $\mathbb{L}$, of length their distance); $|V_Y|_1$ is the sum of the moduli of the coefficients; for $\theta\ge0$ and for any polynomial $F=\sum_YF_Y$ sorted by supports in this way,
\[
\|F\|_\theta\defeq\sup_{x\in\mathbb{L}}\sum_{Y\ni x}e^{\theta\ell(Y)}|F_Y|_1,\qquad\mathsf A\defeq\|V\|_\varsigma\ \text{ for a fixed rate }\varsigma>0,\qquad V_R\defeq\sum_{Y\subset R}V_Y;
\]
\item $\chi_\beta(z)\defeq\prod_{x\in\Lambda}\mathbf 1\{z_x\in\beta K\}$ with $K\subset\R^n$ a closed symmetric convex body, $\bar B(1)\subset K\subset\bar B(c_K)$: the Euclidean ball ($c_K=1$) or Ba{\l}aban's two-component cutoff $K=\{|z^{A}|\le1,\ |z^{\phi}|\le1\}$, $c_K=\sqrt2$, which restricts the photon and the scalar fluctuation at each site separately.
\end{enumerate}
$\kappa_\mu(F_1,\dots,F_k)$ denotes the joint cumulant under a law $\mu$ and $\kappa_\mu(F;k)\defeq\kappa_\mu(F,\dots,F)$ the $k$-th cumulant, the truncated expectation $\langle F^k\rangle^T$. For $X\subset\Lambda$, $P_{XX}$ is the principal block, $G_X\defeq(P_{XX})^{-1}$ and $C\defeq P^{-1}=G_\Lambda$; for $Y\defeq\Lambda\setminus X$ the conditional law of $z_X$ given $z_Y$ under $\mu_P$ is $N(-G_XP_{XY}z_Y,\,G_X)$, by completing the square in the exponent. $\chi^X_\beta\defeq\prod_{x\in X}\mathbf 1\{z_x\in\beta K\}$.

\begin{lemma}[lattice facts]\label{rot:lem:R1lat}
\textup{(i)} $x\mapsto|x-y_0|$ and $x\mapsto\dist(x,S)$ are $1$-Lipschitz. \textup{(ii)} For every $f\ge0$ and $x\in\mathbb{L}$, $\sum_{y\in\mathbb{L}}f(|x-y|)\le\sum_{z\in\Z^3}f(|z|)$. \textup{(iii)} For $x,y\in Y$, $|x-y|\le\ell(Y)$; a shortest connected graph for $Y$ joins any two of its vertices by a path of length at most $\ell(Y)$; and $\ell(Y')\le\ell(Y)$ for $Y'\subset Y$. \textup{(iv)} For $\theta\le\varsigma$ and $r\ge0$ the part of $V$ made of the monomials of length $\ell(Y)\ge r$ has $\|\cdot\|_\theta\le\mathsf Ae^{-(\varsigma-\theta)r}$; in particular $\sum_{Y\ni x,\ \ell(Y)\ge r}|V_Y|_1\le\mathsf Ae^{-\varsigma r}$.
\end{lemma}
\begin{proof}
(i) is the triangle inequality of the quotient metric. (ii) Fix a lift $\tilde x$ and for each $y$ choose the lift $\tilde y$ nearest to $\tilde x$, so that $f(|x-y|)=f(|\tilde x-\tilde y|)$; distinct $y$ have distinct lifts, so the sum is a sub-sum of $\sum_{\tilde z\in\Z^3}f(|\tilde x-\tilde z|)$. (iii) A path in the graph is no longer than the graph, and by the triangle inequality the distance of its endpoints is at most its length; a graph for $Y$ is a graph for $Y'$. (iv) is the definition of the norm, since $e^{\theta\ell}=e^{\varsigma\ell}e^{-(\varsigma-\theta)\ell}$.
\end{proof}

\paragraph{Constants.} Put $\tilde c(\theta)\defeq\sum_{z\in\Z^3}e^{-\theta|z|}$, $\mathsf K\defeq\sum_{z\in\Z^3}|z|e^{-\delta_P|z|/2}$, $c_d\defeq\tilde c(\delta_P/2)$, and
\begin{equation}\label{rot:eq:R1const}
\mu_1\defeq\min\Big\{\frac{\delta_P}2,\frac{\gamma_0}{2c_P\mathsf K}\Big\},\quad g_0\defeq\frac2{\gamma_0},\quad\mathsf M_1\defeq\frac{2c_P}{\gamma_0}\tilde c(\mu_1/2)^2,\quad c_{\mathrm D}\defeq\frac{4c_P}{\gamma_0^2}\tilde c(\mu_1/2)^2,\quad\gamma_K\defeq\min\Big\{\frac12,\frac1{2c_K\mathsf M_1}\Big\},\quad\bar\mu\defeq\min\{\mu_1,\varsigma\}.
\end{equation}
By \Cref{rot:lem:R1lat}(ii) every sum over $\mathbb{L}$ of a function of the distance is bounded by the corresponding sum over $\Z^3$, so these constants serve on every $\mathbb{L}$. They depend on $(\gamma_0,c_P,\delta_P,c_K,\varsigma)$ only, and no lower bound on $c_P$ is used anywhere ($c_P\to0$ gives $\mu_1=\delta_P/2$, $\mathsf M_1\to0$, $\gamma_K=\frac12$).

\begin{lemma}[Combes--Thomas bound]\label{rot:lem:R1ct}
For $X\subset\Lambda$ and $r\ge0$, $\|(P_{XX}+r)^{-1}(x,y)\|\le2(\gamma_0+r)^{-1}e^{-\mu_1|x-y|}$. In particular $\|G_X(x,y)\|\le g_0e^{-\mu_1|x-y|}$ for every $X$, and $\|C(x,y)\|\le g_0e^{-\mu_1|x-y|}$.
\end{lemma}
\begin{proof}
Fix $y_0\in X$ and $f\defeq|\cdot-y_0|$, and conjugate: $P^f\defeq e^{\mu_1f}P_{XX}e^{-\mu_1f}$. Then $(P^f-P_{XX})(x,y)=P(x,y)\big(e^{\mu_1(f(x)-f(y))}-1\big)$ vanishes on the diagonal and, by (R2), \Cref{rot:lem:R1lat}(i) and $e^t-1\le te^t$ for $t\ge0$, has block norm at most $c_Pe^{-\delta_P|x-y|}\mu_1|x-y|e^{\mu_1|x-y|}\le c_P\mu_1|x-y|e^{-\delta_P|x-y|/2}$, since $\mu_1\le\delta_P/2$; by Schur's test and \Cref{rot:lem:R1lat}(ii) its row and column sums are at most $c_P\mu_1\mathsf K\le\gamma_0/2$, so $\|P^f-P_{XX}\|\le\gamma_0/2$. Hence $\Re\langle u,(P^f+r)u\rangle\ge(\gamma_0/2+r)\|u\|^2$ for every $u$ ($P_{XX}\ge\gamma_0$ as a principal block), so $\|(P^f+r)^{-1}\|\le(\gamma_0/2+r)^{-1}\le2(\gamma_0+r)^{-1}$, and $(P_{XX}+r)^{-1}(x,y_0)=e^{-\mu_1f(x)}(P^f+r)^{-1}(x,y_0)e^{\mu_1f(y_0)}$ has norm at most $2(\gamma_0+r)^{-1}e^{-\mu_1|x-y_0|}$.
\end{proof}

\begin{lemma}[conditioning]\label{rot:lem:R1sch}
Let $X\subset\Lambda$ and $Y\defeq\Lambda\setminus X$. Then $0<G_X\le C_{XX}\le\gamma_0^{-1}$, so every component variance of $N(\cdot,G_X)$ and of $\mu_P$ is at most $\gamma_0^{-1}$. The conditional mean $h\defeq-G_XP_{XY}z_Y$ obeys
\[
|h(x)|\le\mathsf M_1e^{-(\mu_1/2)\dist(x,Y)}\|z_Y\|_\infty\qquad(x\in X),
\]
in particular $\|h\|_\infty\le\beta/2$ whenever $\|z_Y\|_\infty\le c_K\gamma_K\beta$. For $x,y\in X$,
\[
\|(C-G_X)(x,y)\|\le c_{\mathrm D}\exp\Big(-\frac{\mu_1}8\big(\dist(x,Y)+\dist(y,Y)\big)-\frac{\mu_1}4|x-y|\Big).
\]
\end{lemma}
\begin{proof}
$C_{XX}=(P_{XX}-P_{XY}P_{YY}^{-1}P_{YX})^{-1}\ge P_{XX}^{-1}=G_X$, because the subtracted operator is non-negative and inversion reverses the order of positive operators; and $C_{XX}\le\|C\|\le\gamma_0^{-1}$. For the mean, \Cref{rot:lem:R1ct} and (R2) (the pairs $(u,v)\in X\times Y$ are off-diagonal) give $|h(x)|\le\sum_{u\in X,v\in Y}g_0e^{-\mu_1|x-u|}c_Pe^{-\delta_P|u-v|}|z_v|$; as $\delta_P\ge2\mu_1$ and $|x-u|+|u-v|\ge|x-v|\ge\dist(x,Y)$, the exponent is at most $-\frac{\mu_1}2\dist(x,Y)-\frac{\mu_1}2(|x-u|+|u-v|)$, and the double sum of the last factor is at most $\tilde c(\mu_1/2)^2$ by \Cref{rot:lem:R1lat}(ii). The particular case is $\mathsf M_1c_K\gamma_K\le\frac12$. For the covariance difference, row $X$ of $PC=\mathbf 1$ reads $P_{XX}C_{XX}+P_{XY}C_{YX}=\mathbf 1$, so $C_{XX}-G_X=-G_XP_{XY}C_{YX}$, with kernel $-\sum_{u\in X,v\in Y}G_X(x,u)P(u,v)C(v,y)$, bounded by $g_0^2c_P\sum_{u,v}e^{-\mu_1L}$ with $L\defeq|x-u|+|u-v|+|v-y|$ (using $\delta_P\ge\mu_1$). Since $L\ge|x-y|$, $L\ge|x-v|\ge\dist(x,Y)$ and $L\ge|v-y|\ge\dist(y,Y)$, write $e^{-\mu_1L}=e^{-\mu_1L/4}e^{-\mu_1L/4}e^{-\mu_1L/2}\le e^{-\frac{\mu_1}8(\dist(x,Y)+\dist(y,Y))}e^{-\frac{\mu_1}4|x-y|}e^{-\frac{\mu_1}2(|x-u|+|u-v|)}$ and sum the last factor over $u,v$.
\end{proof}

\begin{lemma}[decoupling across corridors]\label{rot:lem:R1dec}
Let $\Lambda=\Gamma\sqcup D$ with $D=\bigsqcup_iD_i$ and $\dist(D_i,D_j)\ge w_s$ for $i\ne j$. Let $\mu_P(\cdot\,|\,z_\Gamma)$ be the conditional law of $z_D$ given $z_\Gamma$ and $\nu_{\mathrm{loc}}(\cdot\,|\,z_\Gamma)\defeq\bigotimes_iN(-G_{D_i}P_{D_i\Gamma}z_\Gamma,\,G_{D_i})$ its decoupled version; put $\omega\defeq c_Pc_de^{-\delta_Pw_s/2}$ and suppose $\omega\le\gamma_0/2$. On $\{\|z_D\|_\infty\le\beta,\ \|z_\Gamma\|_\infty\le\beta\}$,
\[
\Big|\log\frac{d\mu_P(\cdot\,|\,z_\Gamma)}{d\nu_{\mathrm{loc}}(\cdot\,|\,z_\Gamma)}\Big|\le\omega_{\mathrm{tot}}\defeq\omega|D|\Big[\frac{\beta^2}2+\frac n{\gamma_0}+\frac{(c_P\tilde c(\delta_P)\beta)^2}{2\gamma_0^2}\Big],
\]
hence $e^{-\omega_{\mathrm{tot}}}\int F\,d\nu_{\mathrm{loc}}\le\int F\,d\mu_P(\cdot\,|\,z_\Gamma)\le e^{\omega_{\mathrm{tot}}}\int F\,d\nu_{\mathrm{loc}}$ for every $F\ge0$ supported in that set. Each factor of $\nu_{\mathrm{loc}}$ is a law of the type of \Cref{rot:lem:R1sch}, with $X=D_i$, $Y=\Lambda\setminus D_i$ and $z_{D_j}=0$ for $j\ne i$.
\end{lemma}
\begin{proof}
Put $u\defeq P_{D\Gamma}z_\Gamma$, $B\defeq\bigoplus_iP_{D_iD_i}$ and $O\defeq P_{DD}-B$ (the blocks $P_{D_iD_j}$, $i\ne j$). The two laws are $N(-P_{DD}^{-1}u,P_{DD}^{-1})$ and $N(-B^{-1}u,B^{-1})$, since $(B^{-1}u)_{D_i}=G_{D_i}P_{D_i\Gamma}z_\Gamma$. The logarithm of the density of $N(-A^{-1}u,A^{-1})$ is $\frac12\log\det(A/2\pi)-\frac12\langle z,Az\rangle-\langle z,u\rangle-\frac12\langle u,A^{-1}u\rangle$; the linear terms agree for $A=P_{DD}$ and $A=B$, so the log-ratio is $\frac12\log\det(P_{DD}B^{-1})-\frac12\langle z_D,Oz_D\rangle-\frac12\langle u,(P_{DD}^{-1}-B^{-1})u\rangle$. Every entry of $O$ is an off-diagonal entry of $P$ between points at distance at least $w_s$, so by (R2) and \Cref{rot:lem:R1lat}(ii) its row and column sums are at most $c_Pe^{-\delta_Pw_s/2}\tilde c(\delta_P/2)=\omega$, and $\|O\|\le\omega$. First term: $|\frac12\langle z_D,Oz_D\rangle|\le\frac12\beta^2\sum_{x\in D}\sum_y\|O(x,y)\|\le\frac12\beta^2|D|\omega$. Second term: $\det(P_{DD}B^{-1})=\det(\mathbf 1+B^{-1/2}OB^{-1/2})$; the symmetric operator $B^{-1/2}OB^{-1/2}$ has $n|D|$ eigenvalues, each of modulus at most $\omega/\gamma_0\le\frac12$ ($B\ge\gamma_0$ as a direct sum of principal blocks), and $|\log(1+t)|\le2|t|$ for $|t|\le\frac12$; so the term is at most $n|D|\omega/\gamma_0$. Third term: $P_{DD}^{-1}-B^{-1}=-P_{DD}^{-1}OB^{-1}$ has norm at most $\omega/\gamma_0^2$, and $|u_x|\le\sum_{v\in\Gamma}c_Pe^{-\delta_P|x-v|}\beta\le c_P\tilde c(\delta_P)\beta$, so $\|u\|^2\le|D|(c_P\tilde c(\delta_P)\beta)^2$. The three bounds add up to $\omega_{\mathrm{tot}}$.
\end{proof}
